\documentclass[10pt]{article}
\usepackage[utf8]{inputenc}
\usepackage[T1]{fontenc}
\usepackage[sc]{mathpazo}
\linespread{1.04}
\usepackage[letterpaper,textwidth=6.7in,textheight=9.5in]{geometry}
\usepackage{booktabs}
\usepackage{tabularx}
\usepackage{caption}
\captionsetup{font=small,labelfont=bf,skip=3pt}
\usepackage{url}
\usepackage[hidelinks]{hyperref}
\usepackage{amsmath}
\usepackage{titlesec}
\titlespacing*{\section}{0pt}{9pt}{4pt}
\titlespacing*{\subsection}{0pt}{6pt}{2pt}
\titleformat{\subsection}[runin]{\normalfont\bfseries}{\thesubsection}{0.6em}{}[.\quad]
\setlength{\parskip}{0pt}
\setlength{\textfloatsep}{8pt}
\setlength{\intextsep}{6pt}
\sloppy
\newcolumntype{Y}{>{\raggedright\arraybackslash}X}
\newcommand{\repo}[1]{\href{https://github.com/Naalu/agent-skill-conflicts/#1}{\texttt{#1}}}

\title{\large\bfseries Checkpoint 2: Research Progress Report\\[3pt]
\normalsize\mdseries Conflicting and competing agent skills in the GitSkills dataset}
\author{\normalsize Chris Reger \and \normalsize Smirthya Somaskantha Iyer}
\date{\normalsize CS 599, Northern Arizona University \\ November 3, 2026}

\begin{document}
\maketitle
\vspace{-14pt}
\noindent Repository: \url{https://github.com/Naalu/agent-skill-conflicts}. Rules, thresholds and seeds are written to
\repo{blob/main/docs/runlog.md} before the code that applies them runs, and corrections are logged there. The paper draft
(introduction through method, result skeleton) is on Overleaf:
\url{https://www.overleaf.com/project/6ae0d951dab4678631de9569}.

\section{Initial research and research questions}

\subsection{Literature}
Every major product shows the model a list of skill names and descriptions and lets it choose; none checks whether
two entries mean the same thing \cite{agentskillsspec2026,claudecodeskills2026,codexskills2026}, and Claude Code's
documentation warns that overlapping descriptions can make it load the wrong skill or miss one. In controlled settings
overlap hurts: one near-synonym competitor per skill cut selection accuracy by 7 to 30\% \cite{li2026singleagent},
routing accuracy falls with the log of library size with errors concentrated among similar skills
\cite{chen2026scaling}, and pass rates fell by up to 21\% at 202 skills because of selection, not context length
\cite{song2026shadowing}. Field prevalence is barely measured: Yan et al.\ \cite{oneskilltoomany2026} found a similar
co-installed skill beside about one in four installed skills across 20,947 repositories with an LLM judge, but exclude
other agents' directories and do not look at drift between copies. Clone research says copies drift and inconsistent
edits cause faults \cite{kim2005genealogies,juergens2009clones}; 53\% of reused skills are never modified after adoption
\cite{gao2026reuse}. Prior overlap detectors use similarity cut-offs from 0.45 to 0.95, mostly unvalidated. An annotated
list of the forty papers we may cite is in \repo{blob/main/docs/references.md}.

\subsection{Dataset}
GitSkills \cite{destefanis2026gitskills} was mined July 9 to 20, 2026. Its 3,797,117 \texttt{SKILL.md} occurrences
collapse to 1,877,981 distinct contents (50.5\% of rows are exact copies of content found elsewhere); one representative
copy per content carries the text and commit history. Front matter parses for 89.2\% of occurrences; commit history
exists for 24.4\% of contents. The dataset holds project-level skills only, so everything we report is a lower bound.

\subsection{Questions}
\textbf{RQ1 (conflicting copies).} How often does a repository give its coding agents different copies of the same
skill, meaning the same name in two agent directories with different content?
\textbf{RQ2 (competing skills).} How often do differently named skills in one catalog compete, meaning a realistic
request could reasonably be answered by either?
\textbf{RQ3 (author practices, secondary).} How often do authors guard against overlap with cross-references, ``do not
use for'' wording, or invocation flags?
The wording is unchanged since Checkpoint 1; the emphasis has moved. The behavioral check once planned under RQ3 is now
an extension if time allows, since Yan et al.\ ran 6,368 Claude Code sessions on it. RQ1 is to our knowledge unmeasured
and needs only names, paths and hashes, so it is an exact count over all eligible repositories. RQ2 replicates the
prevalence question on thirteen times as many repositories, with human labels and a sampling design behind a
reliability gate (Cohen's $\kappa \geq 0.60$ \cite{cohen1960kappa,landis1977kappa}).

\section{Work completed and technical details}

\subsection{Data}
GitSkills July 2026 snapshot, Hugging Face Parquet export (mirror of Zenodo DOI 10.5281/zenodo.21875637), downloaded
October 6 to 8; each \texttt{artifacts} part is read with DuckDB, reduced and deleted (6\,GB to 300\,MB). The unit of
analysis is the repository. A \emph{catalog} is the set of valid skills under one root agent directory
(\texttt{.<agent>/skills/<folder>/SKILL.md}); valid means named exactly \texttt{SKILL.md} with a non-empty front-matter
name and description. This gives 843,271 valid occurrences in 141,458 catalogs across 120,929 repositories, 82,914
catalogs with two or more skills, and 12,875 repositories with skills in two or more directories (the RQ1 population).
Names are keyed on the front-matter field, trimmed and lower-cased, which differs from the folder name in 24.3\% of cases.
Skills with \texttt{disable-model-invocation: true} (2.45\%) are out of RQ2; descriptions are cut at the 1,536-character
product cap before a coder sees them. One e-mail address found in a description was removed from working files, and
row-level files naming repositories stay out of the public repository.

\subsection{Scripts}
Pilot scripts in \repo{tree/main/pilot/collide} (download and reduce, DuckDB tables, RQ1 probes, RQ2 sampling and
all-MiniLM-L6-v2 embeddings, RQ3 counts, an independent recomputation); labelling tools in \repo{tree/main/labelling};
the working-dataset build, validation sampling and GitHub tree checks in \repo{tree/main/src}. Python 3.12, pinned in
\texttt{uv.lock}. Evidence: commit \texttt{ddbd729}, PRs \#16 to \#22, and the \texttt{working-set} branch.

\subsection{Measurement and validation}
RQ1 compares every name present in two directories of one repository by \texttt{file\_sha}, then by whether the
description differs and whether the body length differs by more than 5\% and 50 characters, repeated for the directory
set each product reads. RQ2: two coders label pairs from name and description as C, N or U plus a relation type (D, S,
O) under \repo{blob/main/docs/codebooks/competing-pairs.md}; similarity only stratifies the sample. RQ3: pre-specified
regular expressions and front-matter flags. Checks so far: row counts against the data card; \texttt{verify.py}
reproduced similarity rates on all 20,000 sampled catalogs to within $8\times10^{-7}$; duplicate rates recomputed in SQL
on all 82,914 catalogs; a rebuild after a workspace reset reproduced the pilot result exactly.

\section{Pilot findings}

\begin{table}[t]
\centering\small
\begin{tabular}{lrrrr}
\toprule
Directories compared & Repos, $\geq 2$ dirs & Same name in two & Diverged & \% diverged (95\% CI) \\
\midrule
Any two & 12,875 & 9,875 & 3,591 & 36.4 (35.4--37.3) \\
\quad description differs & & & 1,479 & 15.0 (14.3--15.7) \\
\quad body length differs & & & 1,895 & 19.2 (18.4--20.0) \\
Copilot / VS Code set & 7,428 & 5,777 & 1,906 & 33.0 (31.8--34.2) \\
Cursor set & 8,479 & 6,704 & 2,351 & 35.1 (33.9--36.2) \\
OpenCode set & 6,794 & 5,379 & 1,807 & 33.6 (32.3--34.9) \\
Gemini CLI set & 287 & 205 & 69 & 34.0 (27.5--40.4) \\
\bottomrule
\end{tabular}
\caption{RQ1 at hash level over all 120,929 eligible repositories (\texttt{probes.py}, run log Run 5). Percentages are of
repositories with a same-name skill in two directories; Wilson intervals \cite{wilson1927}. Directory sets per the July
2026 documentation: Copilot \{.github, .claude, .agents\}, Cursor \{.cursor, .agents, .claude, .codex\}, OpenCode
\{.opencode, .claude, .agents\}, Gemini \{.gemini, .agents\}. Claude Code and Codex read one directory each, so they
never see the mirror, only their own copy.}
\label{tab:rq1}
\end{table}

\subsection{RQ1}
Table~\ref{tab:rq1} is the first full-population count. Of the 120,929 repositories, 10.6\% keep skills in two or
more agent directories, 9,875 of those repeat a name across directories (61,159 mirrored names), and in 36.4\% of them at
least one mirrored skill has different content in the two places (15,489 names, 25.3\% of mirrored names). In 15.0\%
the description itself differs, which is the part the agent selects on, and the rate barely changes with the product's
directory set. What this does not say is whether the differences matter: \texttt{Wirasm/PRPs-agentic-eng} differs by one
character of invocation syntax (\texttt{/prp-debug} against \texttt{\$prp-debug}), while
\texttt{wilsonfreitas/awesome-quant} has one copy rewritten in the ``use when the user asks\ldots'' style and one not.
Content has not been diffed yet; a coded sample of diffs is the next step.

\subsection{RQ2 calibration}
Both of us labeled the 180-pair calibration set independently (20 pairs per similarity band; 29 non-English pairs
translated once with DeepL and shared). An LLM reader (C) labeled the same pairs from the codebook alone, for comparison
only (Table~\ref{tab:kappa}).

\begin{table}[h]
\centering\small
\begin{tabular}{llrrrrl}
\toprule
Comparison & Subset & Pairs & \% C, first & \% C, second & $\kappa$ (95\% CI) & Gate \\
\midrule
A vs B & all & 180 & 42 & 56 & 0.61 (0.49--0.72) & met at the point estimate \\
A vs B & English & 151 & & & 0.58 (0.45--0.71) & \\
A vs B & translated & 29 & & & 0.68 (0.40--0.97) & \\
A vs C (LLM) & all & 180 & 42 & 39 & 0.68 (0.57--0.79) & comparison only \\
B vs C (LLM) & all & 180 & 56 & 39 & 0.37 (0.24--0.50) & comparison only \\
\bottomrule
\end{tabular}
\caption{Calibration agreement on competing (C) versus not, from \texttt{labelling/agreement.py}; labels and the
three-way disagreement list are in \repo{tree/main/labelling/calibration-180}. A: lead author; B: co-author.}
\label{tab:kappa}
\end{table}

The gate is met, narrowly: 0.61 against 0.60 with an interval reaching down to 0.49, which we will say wherever the
number appears. The disagreement runs one way: B calls competing 27 pairs that A calls not (C sides with A on 24), the
reverse six times. Nearly all 27 are stages of one workflow (\texttt{rust-reviewer} / \texttt{rust-coder},
\texttt{commit} / \texttt{write-pr}) or one tool on two targets (\texttt{deslop-prose} / \texttt{deslop-code}). A and C read
these as not competing because a realistic request names the stage or target; B read ``same family'' as competing. The
codebook did not pin this down, so version 4 adds a rule for each family, with a worked request, before the main sample
is drawn. The relation types fare worse: on the 70 pairs both called competing, D/S/O agreement is 0.40 ($\kappa$ 0.11),
so types will be reported as exploratory unless a second calibration on types alone reaches 0.60.

\subsection{RQ3}
On the 20,000-catalog pilot sample (192,422 skills): negative-trigger wording in 3.31\% of skills and 10.17\% of
catalogs (6.8\% of catalogs with 2 to 5 skills, 29.0\% with more than 20); a cross-reference to another skill in the same
catalog in 4.76\% of skills and 13.79\% of catalogs (hyphenated names; the pre-specified 6-character rule was inflated by
ordinary-word names); \texttt{disable-model-invocation: true} in 2.45\% of skills, full population. Authors guard against overlap
rarely, mostly in large catalogs; the full-population rerun is pending.

\section{How things are moving}
What went wrong and what we did about it. We wrote the first pilot against a July 31 snapshot date; \texttt{mining\_runs}
showed the last run ended July 20, which moved one figure of the retired-model pilot from 36.0\% to 32.8\%, and July 20
is now used everywhere. We misread Gao et al.'s 58.8\% as the never-modified share. The planned
similarity cut-off for RQ2 failed its pre-registered test (no band reached 80\% competing), so similarity became a
stratifier only; a week lost, but the right design. One repository (\texttt{ssrjkk/claude-skills}) contributed 4,775
generated skills and inflated a dataset-wide percentage by several points before we noticed, which is why we report per
repository. ``One Skill Too Many'' appeared October 8 and took most of the ground planned for the behavioral check; RQ1
became the primary contribution and RQ2 a human-labeled replication, kept mainly because labeling was already under way.
The first pilot readers were shown descriptions without names, which carry much of the signal; hence a third codebook
version. A workspace reset destroyed the pilot database; the rebuild reproduced the result exactly.

The newest problem is the largest. On November 2 we built the working dataset (a pre-registered population ladder:
all repositories, non-forks, project repositories with skill collections set aside, template clusters collapsed) and
drew 200 repositories to validate the project-versus-collection rule. Pulling their file trees from the GitHub API also let
us test whether the dataset holds every byte-identical copy within a repository. It does not. 266 \texttt{SKILL.md} files present on GitHub were missing from the dataset while an identical
copy (same git blob SHA) of each was present elsewhere in the same repository, and the commit history shows 259 of them
existed before the snapshot. GitHub code search often returned one copy of identical content and dropped the rest. Since RQ1 counts identical against diverged mirrors, the dataset's 36.4\% is biased upward and is now
a bound, not the estimate. The fix is to compute RQ1 from each repository's tree at the snapshot commit, which the API gives exactly, on a
stratified sample of about 2,000 project repositories, with the design written before the fetch runs. It is also a
finding about the dataset that we will report.

Other risks: the $\kappa$ interval includes values below the gate (fallback: both-C and either-C bounds with advisor
adjudication); agreement with the LLM reader was near zero on translated pairs; commit dates cover 24\% of contents;
and the week between labels done and analysis frozen is tight.

\section{Remaining work and responsibilities}
\textbf{Done.} Chris: data pipeline, catalog definition, RQ1 count, working dataset and validation sample, RQ2
sampling and embedding code, RQ3 pilot counts, coder A, calibration statistics, repository and run log. Smirthya:
related work and reference list, codebook, coder B, the paper draft (introduction through threats), and Sections 1 and
2 of this report.

\textbf{Remaining.} Codebook v4 with the two rules, logged before sampling (both, Nov 5). Draw and translate the main
sample of about 750 pairs (Chris, Nov 6); label it (both, Nov 15); resolve disagreements and compute agreement (Nov 16).
RQ1: tree-based sample at the snapshot commit, design logged first (Chris, Nov 7); code a random sample of 100 diffs
as substantive or syntactic (both, Nov 12). RQ3 full-population rerun (Chris, Nov 10). Freeze estimates, intervals and figures (Chris, Nov 22).
Results into the draft as numbers land (Smirthya); full draft Dec 1; presentation Dec 3 (Smirthya leads); course paper Dec 9; MSR Dec 18 and 23.

\section{Peer review and response}
We posted the report and pilot analysis to the course channel on November 2 and reviewed two other groups. One review
came back, from Ashley Kalscheur, with four points.

\textbf{The questions are clear, and RQ1 is a specific, measurable contribution.} No change; we agree that RQ1 is the
anchor, and the November 2 finding above is what makes it harder and more interesting than we thought.

\textbf{$\kappa = 0.61$ only narrowly clears 0.60; refine the guidelines and repeat the validation.} We refined: codebook
version 4 adds two rules, each with a worked request, for the two families that account for 27 of the 33
disagreements (stages of one workflow; one tool on two targets), and the rules are frozen before the main sample is
drawn. We did not repeat the 180-pair calibration. A second round would cost both of us a day that the schedule does not
have, and the main sample is itself a larger test of the same thing: agreement on it is reported with its interval, and
if it falls below 0.60 the two-bounds fallback applies. We would rather spend the day on the 750 pairs than on 180 more.

\textbf{Different file contents are not necessarily a real conflict; prioritize the manual diff review.} Agreed, and this
is why the report already leads with ``description changed'' (15.0\%) rather than ``any difference'' (36.4\%), since the
description is the part the agent selects on. The 100-diff coding sample stays on the plan (November 12) and now runs on
the tree-based sample rather than the dataset, for the reason given in Section 4.

\textbf{Three questions is a lot; make RQ3 optional.} Done in substance before the review and now stated plainly: RQ3 is
a set of exact counts that already exist on the pilot sample and need one rerun on the full population; the behavioral
check under it was dropped in October. Nothing else will be added to RQ3.

\small
\bibliographystyle{plain}
\bibliography{references}
\end{document}
